\documentclass[11pt]{article}
\usepackage{fullpage,url,hyperref}
\usepackage{amsmath}
\usepackage{amssymb}

\usepackage[letterpaper,top=1in,bottom=1in,left=1in,right=1in,nohead]{geometry}

\setlength{\parindent}{0in}
\setlength{\parskip}{6pt}

\begin{document}
\thispagestyle{empty}
{\large{\bf ECE 6701 + CS 6501: Optimization for Engineering \hfill Due Fri 10/9}}\\

{\LARGE{\bf Homework 2: Convex Optimization Problems}}
\vspace{0.2\baselineskip}
\hrule

{\bf Instructions:} Submit a single Jupyter notebook (.ipynb) of your work to
Canvas by 11:59pm on the due date. All math must be formatted using LaTeX. {\bf Be sure to show all the work involved in deriving your answers! If you just give
  a final answer without explanation, you may not receive credit for that
  question.}

You may discuss the concepts with your classmates, look on the internet for
help, or ask an LLM for advice, but first try hard to solve problems on your
own, and always write up the answers entirely on your own. Make sure you
understand the concepts and solutions---don't just look up the answers or use AI
to answer it for you. {\bf Cite any sources you used outside of the class
  material (webpages, LLMs, etc.).}

\begin{enumerate}

\item Exercise 4.1 in B\&V. Instead of hand drawing the feasible set, plot it
  using Python (you may find {\tt matplotlib.patches.Polygon} useful for doing
  this.)

\item {\bf Linear Programming.} Consider you need to ship $n$ products from $m$
  warehouses to $n$ destinations. Each warehouse can ship to at most one product
  to a single destination (we can assume $m \geq n$.) The distance from the
  $i$th warehouse to the $j$th destination is given as $d_{ij}$. You want to
  ship all $n$ products with minimal total distance traveled. Set this problem
  up as a linear programming problem, with $x_{ij} = 1$ if the $i$th factory
  ships to the $j$th destination, and $x_{ij} = 0$ otherwise. However, to make
  this a linear programming problem, rather than a discrete optimization, set up
  the optimization variables $x_{ij}$ to be continuous and constrained within
  $[0, 1]$.
  \begin{enumerate}
  \item Write down a formula for the objective function $f(x)$ and the
    inequality / equality constraint functions. This should be in the standard
    form of a linear programming problem. {\bf Note:} Even though the variable
    $x$ is naturally in the form of an $m \times n$ matrix, you can think of it
    as a vector in $\mathbb{R}^{mn}$.

  \item Assuming all of the distances are strictly positive, i.e., $d_{ij} > 0$,
    give an argument for why the optimal solution will actually always be at
    integers, i.e., each $x_{ij}^* = 0$ or $1$.

  \item Use CVXPY to solve the specific problem given by the following distances
    $$
    D = [d_{ij}] = 
    \begin{bmatrix}
      150 & 37 & 88\\
      23 & 42 & 16\\
      105 & 34 & 7\\
      200 & 95 & 12\\
    \end{bmatrix}
    $$

  \item Modify the problem for the case that each factory can ship an unlimited
    number of products. What are the new constraint functions?  Give an analytic
    solution to this problem.

  \item Solve the modified problem using CVXPY for the same distances given in
    part (c), and verify that the answer matches your analytic solution.
  \end{enumerate}

\item Exercise 4.3 in B\&V.

\item Use CVXPY to solve problem 3 above and verify the analytic solution.

\item Exercise 4.5 in B\&V.
  
\item {\bf Quadratic Programming} (Extended version of Exercise 4.22 in B\&V)
  Consider the QCQP
\begin{align*}
\text{minimize} \quad & (1/2)x^T A x + b^Tx + c\\
\text{subject to} \quad & x^T x \leq 0,
\end{align*}
with $A \in S_{++}^n, b \in \mathbb{R}^n,$ and $c \in \mathbb{R}$.
\begin{enumerate}
\item Show that $x^* = -(A + \lambda I)^{-1} b$, where $\lambda = \max\{0,
  \bar{\lambda}\}$, and $\bar{\lambda}$ is the largest solution of the nonlinear
  equation
  $$b^T (A + \lambda I)^{-2} b = 1.$$

\item Consider the relationship of this problem to the group invariance in the
  previous problem (Exercise 4.5). Under what conditions on the parameters
  $A, b, c$ is this QCQP invariant to permutations of $x$? {\bf Hint:} Be sure
  to describe the most general conditions under which permutation invariance
  will hold. The trickiest is the condition on the matrix $A$ - be sure to state
  the condition on $A$ that ensures permutation invariance {\em and} symmetric
  positive-definiteness.

\item Using the result of Exercise 4.5, simplify the above QCQP to a
  one-dimensional optimization problem in the case that permutation invariance
  holds (that is, the conditions you found in part (b) hold.) Write expressions
  for the (univariate) objective function and (univariate) inequality
  constraint.

\item For dimension $n = 8$, pick an explicit set of parameters $A, b, c$ that
  satisfy the permutation invariance conditions in part (b). You {\bf may not}
  pick $A$ to be diagonal, $b = 0$, or $c = 0$, but any other set of parameters
  are okay! (Just to be clear, $A = 0$ is considered a diagonal matrix and is
  thus not allowed.) Now, use CVXPY to solve this problem in the original
  ($n$-dimensional) form. Use CVXPY to solve it again in the reduced
  $1$-dimensional form you derived in part (d) to confirm that the solution is
  the same.
\end{enumerate}

\end{enumerate}

\end{document}
